\subsection{Finite ideal-class norm kernels and split-prime differences}

\begin{statement}{A.92}{Proved}
{Split primes and ramified primes generate the finite class target}
Fix an odd prime $\ell$, a level $n\ge1$, and the actual radical
field $K=K_n$ of A.77.  Write $N=\ell^n$, $L=K(\mu_N)$, and
\[
 G_0=\operatorname{Gal}(L/\Q)
       =(\Z/N\Z)\rtimes(\Z/N\Z)^\times,
 \qquad C=\operatorname{Cl}(K)\{\ell\}.
 \tag{A.92.1}
\]
The group and its natural affine action in (A.92.1) are those
proved in A.77.  For a rational prime $q$, let
\[
 R_q=\operatorname{im}\left(
  \ker\left(\bigoplus_{\mathfrak q\mid q}\Z_\ell[\mathfrak q]
       \xrightarrow{\ \sum f(\mathfrak q/q)\ }\Z_\ell\right)
       \longrightarrow C\right).
 \tag{A.92.2}
\]
Let $C_{\rm split}$ be the subgroup generated by the $R_q$ for
rational primes splitting completely in $L$, and let
$C^0=\operatorname{im}(W_n\to C)$ as in A.86.  Then
\[
 C^0=C_{\rm split}
       +\sum_{q\ {
m ramified\ in}\ L/\Q}R_q.
 \tag{A.92.3}
\]
In defining $C_{\rm split}$ one may omit any prescribed finite
set of rational primes.  Since $C$ is finite, a finite set of
the remaining split primes already generates $C_{\rm split}$.

For a ramified $q\ne\ell$, the map from A.89 identifies
$W_{\infty,q}$ with the local ideal-norm kernel in (A.92.2)
at this level.  For $\ell=3$, the remaining coefficient-place
kernel is zero.  Consequently, with the actual global inverse
system of A.86, one also has the finite-level equality
\[
 C_n^0=C_{n,\rm split}
       +\operatorname{im}\left(
          \varprojlim_jW_j\longrightarrow W_n\longrightarrow C_n^0
                            \right)\qquad(\ell=3).
 \tag{A.92.4}
\]
This is an equality at a fixed level.  It does not assert that
the split-prime generators have compatible representatives in
$\varprojlim W_j$, or that the particular class family of A.86
lies in its image.
\end{statement}
\begin{proof}
All class-field maps used in the argument are specified below.
Let $H_K/K$ be the primary Hilbert class field, with the arithmetic
Artin identification
$\operatorname{Gal}(H_K/K)=\operatorname{Cl}(K)\{\ell\}=C$.
This is the primary part of the Hilbert class field correspondence
\cite[V, Theorems 3.5--3.6 and Example 3.9]{MilneCFT}.
Let $F$ be the compositum of $L$ and the finitely many
$\Q$-conjugates of $H_K$.  It is a finite Galois extension of
$\Q$.  Put
\[
 G=\operatorname{Gal}(F/\Q),\qquad
 H=\operatorname{Gal}(F/K),\qquad
 V=\operatorname{Gal}(F/L),\qquad U=(\Z/N\Z)^\times.
 \tag{A.92.5}
\]
Each base change of a conjugate Hilbert class field is an
unramified abelian $\ell$-extension of $L$.  Therefore $V$ is
an abelian $\ell$-group and $F/L$ is unramified at every finite
prime.  Restriction gives the exact quotient $G/V=G_0$ and
the root stabilizer $H/V=U$.  Restriction also gives a
surjection $H\to C$.

We first prove directly the extension property of characters
needed here.  Give $\Q_\ell/\Z_\ell$ its trivial finite-group
action.  For odd $\ell$,
\[
                  H^2(G_0,\Q_\ell/\Z_\ell)=0.
 \tag{A.92.6}
\]
Indeed, apply the Hochschild--Serre spectral sequence to
$1\to\Z/N\Z\to G_0\to U\to1$.
The groups $U$ and $\Z/N\Z$ are cyclic, and degree-two
cohomology of a finite cyclic group with this divisible
trivial module is zero.  The remaining term of total degree
two is
$H^1(U,\operatorname{Hom}(\Z/N\Z,\Q_\ell/\Z_\ell))$.
Choose a generator $r$ of $U$.  On this character module it
acts by $r^{-1}$.  As $r-1$ is invertible modulo $\ell$, the
operator $r^{-1}-1$ is an automorphism.  The cyclic group
cochain formula therefore gives zero for this $H^1$ as well.
All terms of total degree two vanish, proving (A.92.6).
The same commutator calculation gives
$G_0^{\rm ab}=U$: conjugation by $r$ sends a translation
$t$ to $rt$, and $(r-1)(\Z/N\Z)=\Z/N\Z$.

Let $\psi:C\to\Q_\ell/\Z_\ell$ be a character and also denote
its pullback to $H$ by $\psi$.
We claim that $\psi$ annihilates $C_{\rm split}$ exactly when
its restriction to $V$ is invariant under $G$.
For any $\sigma\in V$, Chebotarev gives rational primes
outside any prescribed finite set whose Frobenius in $F$
has conjugacy class that of $\sigma$.
Here we use the finite-conjugacy-class form
\cite[Theorem 8.31]{AgentMilneANT}.
Choose a prime of $F$ for which Frobenius is $\sigma$ itself.
These rational primes split completely in $L$.
Their primes in $K$ are indexed by $G/H$, all have residue
degree one, and the Artin value at the branch indexed by
$gH$ is the restriction to $H_K$ of $g^{-1}\sigma g\in H$.
It follows that the character values on these prime classes
are exactly
\[
                         \psi(g^{-1}\sigma g),\qquad gH\in G/H.
 \tag{A.92.7}
\]
The kernel in (A.92.2) for such a prime is generated over
$\Z_\ell$ by differences of its prime basis vectors.
Thus annihilation of all these differences is exactly
constancy of (A.92.7), for every $\sigma\in V$.
Since $\psi$ is a character of $H$, this constancy is
equivalent to invariance of $\psi|_V$ under all of $G$.
The same proof holds after omitting any finite set of primes.

Suppose now that $\psi$ annihilates $C_{\rm split}$.
The five-term sequence for $1\to V\to G\to G_0\to1$,
together with (A.92.6), supplies a character
$\chi:G\to\Q_\ell/\Z_\ell$ whose restriction to $V$ is
$\psi|_V$.
The difference $\psi-\chi|_H$ factors through $H/V=U$.
The quotient map $G\to G_0\to U$ restricts on $H/V$ to
the identity, so this difference itself extends to $G$.
Adding that extension to $\chi$, we obtain an exact extension
\[
                              \chi|_H=\psi.
 \tag{A.92.8}
\]

Take any rational prime $q$ unramified in $L$.
It is also unramified in $F$, because $F/L$ is unramified.
Choose its Frobenius $\sigma\in G$.
Its cycles on $G/H$ index the primes
$\mathfrak q_1,\ldots,\mathfrak q_t$ of $K$ above $q$;
their lengths are the actual residue degrees $f_1,\ldots,f_t$.
For a representative $g_iH$ of the $i$th cycle, the
corresponding Frobenius over $K$ is
$g_i^{-1}\sigma^{f_i}g_i\in H$.
Using (A.92.8) and that $\chi$ is a character gives
\[
 \psi([\mathfrak q_i])
       =\chi(g_i^{-1}\sigma^{f_i}g_i)
       =f_i\chi(\sigma).
 \tag{A.92.9}
\]
Consequently $\psi$ vanishes on $R_q$, since every vector
in its source satisfies $\sum_i f_i t_i=0$.
Characters of the finite abelian $\ell$-group $C$ separate
subgroups.  Applying this to all characters annihilating
$C_{\rm split}$ proves
$R_q\subset C_{\rm split}$ for every $q$ unramified in $L$.

The completed group $W_n$ is specified by its local vectors,
each of which is in the kernel in (A.92.2).
Let $\ell^a$ annihilate $C$.
The convergence condition in (A.86.6) says that all but
finitely many local vectors have every entry divisible by
$\ell^a$.  Such a local vector is $\ell^a$ times another
vector in the same kernel, because $\Z_\ell$ is torsion-free.
Their completed sum is also $\ell^a$ times a completed
divisor, so it has zero class in $C$.
Thus $C^0$ is the algebraic sum of the $R_q$.
The preceding inclusion for unramified primes proves
(A.92.3).  The character argument with a finite excluded
set proves that its omission does not change
$C_{\rm split}$; finite generation then follows from the
finiteness of $C$.

Finally let $q\ne\ell$ ramify in $L$.
The cyclotomic extension is unramified at $q$, and the
tame Kummer inertia on the root labels is translation by
$m_q=v_q(A)\pmod{\ell^n}$.
Ramification is therefore equivalent to
$s_q=v_\ell(m_q)<n$.
The stable description (A.89.3) applies already at level
$n$, and identifies its local inverse kernel with the
kernel in (A.92.2).  A finite family of these local lifts
belongs to the global inverse system: it has finite support,
so the exact convergence condition of (A.89.4) holds.
When $\ell=3$, the Eisenstein calculation in A.89 gives
$W_{n,3}=0$ at every level.  All ramified-prime terms in
(A.92.3) therefore lie in the image displayed in (A.92.4).
The reverse inclusion is contained in $C_n^0$ by definition,
which proves that equality and its stated finite-level scope.
\end{proof}

The relative transition problem remains explicit.  If an ideal
class at level $n+1$ has zero class norm, an ideal representing
it may have ideal norm $\operatorname{div}(a)$ rather than zero.
Making that ideal norm literally zero by a principal correction
requires $a$, up to a global unit, to be a field norm.
Formula (A.92.3) does not establish this relative assertion.
It also does not supply the levelwise convergence condition
needed to sum infinitely many split-prime corrections.

\noindent The exact finite certificate
\path{code/norm_split_generation_certificate.py}, with retained output
\path{code/norm_split_generation_certificate_output.txt}, checks the
cyclic character differential, the affine commutator and stabilizer
maps, all Frobenius-cycle identities in the listed finite groups,
and the split-character criterion in explicit permutation-module
extensions.  The number-field class-field correspondence and
Chebotarev steps are proved and cited above.
